% Unit 092 terjemahan Bahasa Indonesia dari chapter7.tex baris 892--1129.
% Sumber: Methods of Algebra, Volume 2, commit
% 9a5803ff77dd3257484cb177f851a73770a59dd3 (CC BY 4.0).
% stable-unit-id: o014.aljabr2.chapter7.coalgebra-to-hopf-a
% Cakupan: Bagian 7.5 dari pembukaan hingga bukti struktur monoidal kategori
% modul atas bialjabar; dilanjutkan oleh Unit 093 pada baris 1130.

% segment-id: o014.li2.chapter7.coalgebra-to-hopf-a.h001
\section{Dari Koaljabar ke Aljabar Hopf}\label{sec:bialgebra}
\hypertarget{o014.aljabr2.chapter7.coalgebra-to-hopf-a}{ }
% segment-id: o014.li2.chapter7.coalgebra-to-hopf-a.p001
Tujuan pertama bagian ini ialah mendualkan Definisi~\ref{def:algebra-monoidal}
tentang aljabar. Untuk itu, mula-mula amati bahwa definisi kategori monoidal
bersifat swadual. Misalkan $\mathcal{V}$ kategori monoidal. Pada
$\mathcal{V}^{\opp}$, kita tetap dapat memakai bifunktor $\otimes$ dan objek
satuan $\munit$ semula, tetapi kendala asosiativitas $a$ dan
$\iota:\munit\otimes\munit\rightiso\munit$ dalam
\cite[Definisi 3.1.1]{Li1} diganti dengan inversnya. Dengan demikian,
$\mathcal{V}^{\opp}$ menjadi kategori monoidal.

% segment-id: o014.li2.chapter7.coalgebra-to-hopf-a.p002
Serupa dengan itu, jika $\mathcal{V}^{\opp}$ memiliki struktur anyaman $c$
seperti dalam \cite[Definisi 3.3.1]{Li1}, mengambil invers memberikan
struktur anyaman yang bersesuaian pada $\mathcal{V}^{\opp}$. Dalam
korespondensi ini, $\mathcal{V}$ adalah kategori monoidal simetris jika dan
hanya jika $\mathcal{V}^{\opp}$ juga demikian.

% segment-id: o014.li2.chapter7.coalgebra-to-hopf-a.d001
\begin{definition}[Koaljabar]\label{def:cogebra-monoidal}
	\index{koaljabar@koaljabar (coalgebra)}
	Misalkan $\mathcal{V}$ kategori monoidal. Aljabar di atas
	$\mathcal{V}^{\opp}$ disebut \emph{koaljabar} di atas $\mathcal{V}$.
	Dengan kata lain, koaljabar adalah data $(C,\Delta,\epsilon)$, dengan
	% segment-id: o014.li2.chapter7.coalgebra-to-hopf-a.q001
	\[ C \in \Obj(\mathcal{V}), \quad \Delta = \Delta_C: C \to C \otimes C, \quad \epsilon = \epsilon_C: C \to \munit, \]
	yang memenuhi diagram komutatif dual dari diagram dalam
	Definisi~\ref{def:algebra-monoidal}:
	% segment-id: o014.li2.chapter7.coalgebra-to-hopf-a.q002
	\[\begin{tikzcd}
			\munit \otimes C & C \otimes C \arrow[r, "{\identity \otimes \epsilon}"] \arrow[l, "{\epsilon \otimes \identity}"'] & C \otimes \munit \\
			& C \arrow[lu, "\sim"' sloped] \arrow[ru, "\sim"' sloped] \arrow[u, "\Delta"'] &
		\end{tikzcd} \quad
	\begin{tikzcd}[column sep=tiny]
		(C \otimes C) \otimes C & & C \otimes (C \otimes C) \arrow[ll, "\sim"] \\
		C \otimes C \arrow[u, "{\Delta \otimes \identity}"] & C \arrow[l, "\Delta"] \arrow[r, "\Delta"'] & C \otimes C . \arrow[u, "{\identity \otimes \Delta}"']
	\end{tikzcd}\]

	Lazimnya, $\Delta$ disebut \emph{komultiplikasi} pada $C$, sedangkan
	$\epsilon$ disebut \emph{kounit} pada $C$.

	Morfisme dari koaljabar $(C,\Delta,\epsilon)$ ke
	$(C',\Delta',\epsilon')$ didefinisikan secara dual terhadap
	Definisi~\ref{def:algebra-monoidal}: morfisme itu adalah
	$\phi:C\to C'$ yang membuat diagram berikut komutatif:
	% segment-id: o014.li2.chapter7.coalgebra-to-hopf-a.q003
	\[\begin{tikzcd}
		\munit & C \arrow[d, "\phi"] \arrow[l, "\epsilon"']  \\
		& C' \arrow[lu, "{\epsilon'}"]
	\end{tikzcd} \quad \begin{tikzcd}
		C \otimes C \arrow[r, "{\phi \otimes \phi}"] & C' \otimes C' \\
		C \arrow[r, "\phi"'] \arrow[u, "\Delta"] & C' . \arrow[u, "{\Delta'}"']
	\end{tikzcd}\]

	Jika $\mathcal{V}$ dilengkapi struktur anyaman $c$ dan $c\Delta=\Delta$,
	maka $(C,\Delta,\epsilon)$ disebut \emph{kokomutatif}.
\end{definition}

% segment-id: o014.li2.chapter7.coalgebra-to-hopf-a.p003
Definisi berikut tidak lain adalah dual dari
Definisi~\ref{def:module-algebra}.

% segment-id: o014.li2.chapter7.coalgebra-to-hopf-a.d002
\begin{definition}[Komodul]\label{def:comodule-cogebra}
	\index{komodul@komodul (comodule)}
	Misalkan $(C,\Delta,\epsilon)$ koaljabar di atas $\mathcal{V}$. Komodul-$C$
	kiri didefinisikan sebagai data $(M,\rho)$, dengan
	% segment-id: o014.li2.chapter7.coalgebra-to-hopf-a.q004
	\[ M \in \Obj(\mathcal{V}), \quad \rho = \rho_M : M \to C \otimes M, \]
	dan $\rho$ dapat dipandang sebagai komultiplikasi skalar pada komodul kiri.
	Syaratnya ialah komutativitas diagram berikut:
	% segment-id: o014.li2.chapter7.coalgebra-to-hopf-a.q005
	\[\begin{tikzcd}
		\munit \otimes M & C \otimes M \arrow[l, "{\epsilon \otimes \identity}"'] \\
		& M \arrow[u, "{\rho}"'] \arrow[lu, "\sim" sloped]
	\end{tikzcd} \quad \begin{tikzcd}[column sep=small]
		(C \otimes C) \otimes M & & C \otimes (C \otimes M) \arrow[ll, "\sim"'] \\
		C \otimes M \arrow[u, "{\Delta \otimes \identity}"] & M \arrow[l, "{\rho}"] \arrow[r, "{\rho}"'] & C \otimes M. \arrow[u, "{\identity \otimes \rho}"']
	\end{tikzcd}\]
	Homomorfisme dari komodul kiri $(M,\rho)$ ke $(M',\rho')$ adalah morfisme
	$f:M\to M'$ yang memenuhi
	$(\identity\otimes f)\rho=\rho'f$. Komodul-$C$ kanan, bikomodul, dan
	homomorfisme di antaranya didefinisikan dengan cara serupa.
\end{definition}

% segment-id: o014.li2.chapter7.coalgebra-to-hopf-a.p004
Contoh awal yang penting ialah $C$ sendiri: dengan
$\rho_C:=\Delta:C\to C\otimes C$, objek ini menjadi komodul kiri sekaligus
komodul kanan.

% segment-id: o014.li2.chapter7.coalgebra-to-hopf-a.p005
Walaupun definisi komodul tampak berkebalikan dengan kebiasaan, dalam praktik
komodul belum tentu lebih rumit daripada modul. Sebagai contoh, ambil
$\mathcal{V}=\Bbbk\dcate{Mod}$ dengan $\Bbbk$ gelanggang komutatif. Misalkan
$M$ modul-$\Bbbk$ bebas dengan basis $(v_i)_{i\in I}$. Memberikan
$\rho:M\to M\dotimes{\Bbbk}C$ sama artinya dengan memberikan keluarga unsur
$(t_{ij})_{(i,j)\in I^2}$ dari $C$ sedemikian sehingga, untuk setiap $i\in I$,
% segment-id: o014.li2.chapter7.coalgebra-to-hopf-a.q006
\begin{equation*}
	\rho(v_i) = \sum_{j \in I} v_j \otimes t_{ji} \quad \text{(jumlah hingga)},
\end{equation*}
dan syarat komodul dinyatakan, untuk setiap $i$, sebagai
% segment-id: o014.li2.chapter7.coalgebra-to-hopf-a.q007
\begin{equation*}
	v_i = \sum_j \epsilon(t_{ji}) v_j, \quad \sum_j v_j \otimes \Delta(t_{ji}) = \sum_{j, k} v_k \otimes t_{kj}  \otimes t_{ji}.
\end{equation*}
Dengan mengganti nama indeks secara tepat, syarat ini selanjutnya
disederhanakan, untuk setiap $i,j$, menjadi
% segment-id: o014.li2.chapter7.coalgebra-to-hopf-a.q008
\begin{equation}\label{eqn:comodule-matrix}
	\begin{aligned}
		\epsilon(t_{ji}) & = \delta_{j, i}, \\
		\Delta(t_{ji}) & = \sum_k t_{jk} \otimes t_{ki};
	\end{aligned}
\end{equation}
di sini $\delta_{j,i}$ adalah simbol delta Kronecker, yang didefinisikan
sebagai $1$ bila $j=i$ dan sebagai $0$ bila tidak.
\index{Kronecker-delta@simbol $\delta$ Kronecker (Kronecker delta)}
\index[sym1]{delta-ij@$\delta_{ij}$}

% segment-id: o014.li2.chapter7.coalgebra-to-hopf-a.p006
Melanjutkan pembahasan di atas, tuliskan pemetaan $\Bbbk$-linear
$f:M\to M'$ antara komodul kanan terhadap basis yang dipilih sebagai
$f(v_i)=\sum_jr_{ji}v'_j$, yakni sebagai matriks dengan
$r_{ji}\in\Bbbk$. Dengan cara yang sama dapat dibuktikan bahwa $f$ adalah
homomorfisme komodul jika dan hanya jika, untuk semua $i,k$,
% segment-id: o014.li2.chapter7.coalgebra-to-hopf-a.q009
\begin{equation*}\begin{gathered}
	\sum_j r_{kj} t_{ji} = \sum_j r_{ji} t'_{kj},
\end{gathered}\end{equation*}
dengan $\Bbbk$ beraksi dari kiri pada $C$ melalui perkalian skalar. Versi
untuk komodul kiri tentu serupa.

% segment-id: o014.li2.chapter7.coalgebra-to-hopf-a.p007
Kita akan melanjutkan pembahasan beberapa aspek teori komodul dalam
\S\ref{sec:Morita}.

% segment-id: o014.li2.chapter7.coalgebra-to-hopf-a.r001
\begin{remark}\label{rem:coalgebra-convolution}
	\index{konvolusi@konvolusi (convolution)}
	Jika $(C,\Delta,\epsilon)$ koaljabar di atas $\mathcal{V}$ dan
	$(A,\mu,\eta)$ aljabar di atas $\mathcal{V}$, maka $\Hom(C,A)$ dilengkapi
	operasi biner $\star$ yang disebut \emph{konvolusi}, sebagai berikut:
	% segment-id: o014.li2.chapter7.coalgebra-to-hopf-a.q010
	\[ f \star g := \mu (f \otimes g) \Delta, \quad f, g \in \Hom(C, A). \]
	Dari sifat-sifat $\Delta$ dan $\mu$, mudah dilihat bahwa $\star$ asosiatif,
	dan diagram komutatif berikut cukup untuk menunjukkan bahwa
	$(\Hom(C,A),\star)$ adalah monoid dengan satuan
	$\eta\epsilon\in\Hom(C,A)$:
	% segment-id: o014.li2.chapter7.coalgebra-to-hopf-a.q011
	\[\begin{tikzcd}[column sep=large]
		A \otimes \munit \arrow[r, "{\mu(\identity \otimes \eta)}"] & A & \munit \otimes A \arrow[l, "{\mu(\eta \otimes \identity)}"'] \\
		C \otimes \munit \arrow[r] \arrow[u, "{f \otimes \identity}"] & C \arrow[u, "f"] & \munit \otimes C \arrow[l] \arrow[u, "{\identity \otimes f}"'] \\
		C \otimes C \arrow[u, "{\identity \otimes \epsilon}"] & C \arrow[l, "\Delta"] \arrow[r, "\Delta"'] \arrow[equal, u] & C \otimes C . \arrow[u, "{\epsilon \otimes \identity}"']
	\end{tikzcd}\]
	Argumen standar menunjukkan bahwa jika $\varphi:C'\to C$ morfisme
	antarkoaljabar dan $\psi:A\to A'$ morfisme antaraljabar, maka
	% segment-id: o014.li2.chapter7.coalgebra-to-hopf-a.q012
	\begin{equation}\label{eqn:convolution-functorial}
		\Hom(C, A) \to \Hom(C', A'), \quad f \mapsto \psi f \varphi
	\end{equation}
	adalah homomorfisme antara monoid-monoid konvolusi.
\end{remark}

% segment-id: o014.li2.chapter7.coalgebra-to-hopf-a.p008
Mulai sekarang, misalkan $\mathcal{V}$ kategori monoidal beranyam, dan
tuliskan struktur anyamannya seperti biasa dalam bentuk
% segment-id: o014.li2.chapter7.coalgebra-to-hopf-a.q013
\[ c(X, Y): X \otimes Y \to Y \otimes X. \]
Menurut \S\ref{sec:alg-in-monoidal-cat}, kategori aljabar
$\cate{Alg}(\mathcal{V})$ (atau kategori koaljabar
$\cate{Alg}(\mathcal{V}^{\opp})^{\opp}$) mempunyai struktur monoidal dengan
satuan $\munit$. Karena itu, kita dapat meninjau koaljabar di dalam kategori
aljabar (atau aljabar di dalam kategori koaljabar). Pengamatan sederhana
namun menarik ialah bahwa keduanya mengarah pada data yang sama
$(A,\mu,\eta,\Delta,\epsilon)$, dengan $A\in\Obj(\mathcal{V})$. Syarat yang
diperlukan bagi morfisme
$A\xrightarrow{\Delta}A\otimes A\xrightarrow{\mu}A$ dan
$\munit\xrightarrow{\eta}A\xrightarrow{\epsilon}\munit$ dalam
$\mathcal{V}$ adalah:
% segment-id: o014.li2.chapter7.coalgebra-to-hopf-a.l001
\begin{enumerate}[(i)]
	\item $(A,\mu,\eta)$ adalah aljabar;
	\item $(A,\Delta,\epsilon)$ adalah koaljabar;
	\item komultiplikasi $\Delta:A\to A\otimes A$ dan kounit
	$\epsilon:A\to\munit$ adalah morfisme antaraljabar;
	\item perkalian $\mu:A\otimes A\to A$ dan satuan
	$\eta:\munit\to A$ adalah morfisme antarkoaljabar.
\end{enumerate}
Setelah (i) dan (ii) diasumsikan, (iii) dan (iv) ekuivalen, sehingga cukup
memeriksa salah satunya. Keduanya menyatakan kompatibilitas antara
$(\Delta,\epsilon)$ dan $(\mu,\eta)$ melalui empat diagram komutatif untuk
empat kemungkinan pasangan, atau dalam rumus,
% segment-id: o014.li2.chapter7.coalgebra-to-hopf-a.q014
\begin{equation*}\begin{gathered}
	\Delta \mu = (\mu \otimes \mu)(\identity_A \otimes c(A, A) \otimes \identity_A) (\Delta \otimes \Delta), \\
	\Delta \eta = \eta \otimes \eta, \quad \epsilon \mu = \epsilon \otimes \epsilon, \quad \epsilon \eta = \identity_{\munit}.
\end{gathered}\end{equation*}

% segment-id: o014.li2.chapter7.coalgebra-to-hopf-a.d003
\begin{definition}[Bialjabar]
	\index{bialjabar@bialjabar (bialgebra)}
	Misalkan $\mathcal{V}$ kategori monoidal beranyam. Jika $A$ adalah
	aljabar dalam kategori koaljabar di atas $\mathcal{V}$, atau secara
	ekuivalen koaljabar dalam kategori aljabar di atas $\mathcal{V}$, maka $A$
	disebut \emph{bialjabar} di atas $\mathcal{V}$. Secara konkret, bialjabar
	terdiri atas data $(A,\mu,\eta,\Delta,\epsilon)$.

	Jika $A$ dan $A'$ keduanya bialjabar dan $\phi:A\to A'$ sekaligus
	merupakan morfisme antaraljabar dan morfisme antarkoaljabar, maka $\phi$
	disebut morfisme dari bialjabar $A$ ke $A'$.
\end{definition}

% segment-id: o014.li2.chapter7.coalgebra-to-hopf-a.p009
Seluruh data suatu bialjabar sering disingkat hanya dengan $A$.

% segment-id: o014.li2.chapter7.coalgebra-to-hopf-a.p010
Jika $\Bbbk$ gelanggang komutatif dan
$\mathcal{V}=\Bbbk\dcate{Mod}$, maka aljabar, koaljabar, dan bialjabar di atas
$\mathcal{V}$ juga masing-masing disebut aljabar-$\Bbbk$,
koaljabar-$\Bbbk$, dan bialjabar-$\Bbbk$, dan seterusnya.

% segment-id: o014.li2.chapter7.coalgebra-to-hopf-a.eg001
\begin{example}\label{eg:monoid-bialgebra}
	Misalkan $M$ monoid dan $\Bbbk$ gelanggang komutatif. Bentuk aljabar monoid
	$\Bbbk[M]$ di atas $\Bbbk$; lihat \cite[Definisi 5.6.1]{Li1}. Definisikan
	homomorfisme modul-$\Bbbk$
	% segment-id: o014.li2.chapter7.coalgebra-to-hopf-a.q015
	\begin{align*}
		\Delta: \Bbbk[M] & \to \Bbbk[M] \dotimes{\Bbbk} \Bbbk[M] \\
		m & \mapsto m \otimes m, \quad m \in M
	\end{align*}
	dan $\epsilon:\Bbbk[M]\to\Bbbk$ dengan $\epsilon$ memetakan setiap
	$m\in M$ ke $1$. Mari kita periksa bahwa data ini memberikan
	bialjabar-$\Bbbk$.

	Pertama, $(\Bbbk[M],\Delta,\epsilon)$ adalah koaljabar, sebab untuk setiap
	$m\in M$ berlaku
	% segment-id: o014.li2.chapter7.coalgebra-to-hopf-a.q016
	\begin{gather*}
		(\identity \otimes \epsilon)(\Delta(m)) = m = (\epsilon \otimes \identity)(\Delta(m)), \\
		(\Delta \otimes \identity)(\Delta(m)) = m \otimes m \otimes m = (\identity \otimes \Delta)(\Delta(m)).
	\end{gather*}
	Perhatikan bahwa $\Bbbk[M]$ selalu kokomutatif, tetapi $\Bbbk[M]$
	komutatif jika dan hanya jika $M$ komutatif.

	Kedua, periksa bahwa $\epsilon$ dan $\Delta$ keduanya homomorfisme
	aljabar-$\Bbbk$. Cukup diperiksa, untuk setiap
	$m,m'\in M$, bahwa
	% segment-id: o014.li2.chapter7.coalgebra-to-hopf-a.q017
	\begin{gather*}
		\epsilon(mm') = 1 = \epsilon(m)\epsilon(m'), \quad \epsilon(1) = 1, \\
		\Delta(mm') = mm' \otimes mm' = (m \otimes m)(m' \otimes m') = \Delta(m) \Delta(m'), \\
		\Delta(1) = 1 \otimes 1.
	\end{gather*}
\end{example}

% segment-id: o014.li2.chapter7.coalgebra-to-hopf-a.p011
Untuk aljabar $A$ di atas $\mathcal{V}$,
Definisi~\ref{def:module-algebra} menjelaskan modul-$A$ kiri. Jika $A$
bialjabar, kita dapat melangkah lebih jauh sebagai berikut.
% segment-id: o014.li2.chapter7.coalgebra-to-hopf-a.l002
\begin{enumerate}
	\item Pada objek satuan $\munit$, definisikan struktur modul-$A$ kiri
	melalui morfisme $\mu_{\munit}:A\otimes\munit\to\munit$ yang merupakan
	komposisi
	$A\otimes\munit\rightiso A\xrightarrow{\epsilon}\munit$.
	\item Untuk modul-$A$ kiri $M_i$, tuliskan morfisme perkalian skalar yang
	bersesuaian sebagai $\mu_i$ ($i=1,2$). Definisikan morfisme
	$\mu_{M_1\otimes M_2}:A\otimes(M_1\otimes M_2)\to M_1\otimes M_2$ sebagai
	komposisi
	% segment-id: o014.li2.chapter7.coalgebra-to-hopf-a.q018
	\begin{multline*}
		A \otimes M_1 \otimes M_2 \xrightarrow{\Delta \otimes \identity_{M_1 \otimes M_2}} A \otimes A \otimes M_1 \otimes M_2 \\
		\xrightarrow{\identity_A \otimes c(A, M_1) \otimes \identity_{M_2}} A \otimes M_1 \otimes A \otimes M_2 \xrightarrow{\mu_1 \otimes \mu_2} M_1 \otimes M_2.
	\end{multline*}
	Di sini, untuk menyederhanakan notasi, tanda kurung operasi $\otimes$ dan
	kendala asosiativitas dihilangkan; dengan kata lain, $\mathcal{V}$
	diperlakukan sebagai kategori monoidal ketat.
\end{enumerate}

% segment-id: o014.li2.chapter7.coalgebra-to-hopf-a.p012
Definisi ini dimaksudkan untuk memberi objek $\munit$ dan
$M_1\otimes M_2$ dari $\mathcal{V}$ struktur modul-$A$ kiri. Kasus modul-$A$
kanan tentu serupa.

% segment-id: o014.li2.chapter7.coalgebra-to-hopf-a.pr001
\begin{proposition}\label{prop:bialgebra-Mod-monoidal}
	\index[sym1]{Comod@$\cate{Comod}$}
	Misalkan $A$ bialjabar di atas kategori monoidal beranyam $\mathcal{V}$.
	Definisi di atas membuat kategori modul-$A$ kiri $A\dcate{Mod}$ menjadi
	kategori monoidal dengan satuan $(\munit,\mu_{\munit})$. Hal yang sama
	berlaku bagi kategori modul-$A$ kanan $\cated{Mod}A$.

	Secara dual, perkalian dan satuan $A$ membuat kategori komodul-$A$ kiri
	$A\dcate{Comod}$ dan kategori komodul-$A$ kanan $\cated{Comod}A$ menjadi
	kategori monoidal.

	Untuk keempat kategori di atas, funktor pelupa ke $\mathcal{V}$ selalu
	merupakan funktor monoidal.
\end{proposition}
% segment-id: o014.li2.chapter7.coalgebra-to-hopf-a.pf001
\begin{proof}
	Berdasarkan dualitas, cukup dibuktikan bagian pertama. Pemeriksaannya panjang
	tetapi tidak sulit; garis besarnya sebagai berikut. Pertama-tama harus
	diperiksa bahwa perkalian skalar memenuhi diagram komutatif dalam
	Definisi~\ref{def:module-algebra}. Untuk $(\munit,\mu_{\munit})$,
	sifat-sifat tersebut mengikuti fakta bahwa
	$\epsilon:A\to\munit$ adalah morfisme aljabar. Untuk struktur modul pada
	$M_1\otimes M_2$, diagram tentang perkalian skalar oleh satuan mengikuti
	fakta bahwa $\Delta:A\to A\otimes A$ adalah morfisme aljabar,
	khususnya bahwa ia kompatibel dengan $\eta$ dan $\eta\otimes\eta$.

	Untuk asosiativitas perkalian skalar, setelah definisinya diuraikan,
	persoalannya ialah membuktikan kesamaan dua pemetaan dari
	$A\otimes A\otimes M_1\otimes M_2$ ke $M_1\otimes M_2$:
	% segment-id: o014.li2.chapter7.coalgebra-to-hopf-a.l003
	\begin{enumerate}[(a)]
		\item mula-mula terapkan
		$\Delta\otimes\Delta\otimes\identity_{M_1}\otimes\identity_{M_2}$,
		lalu, dengan memperhitungkan struktur anyaman, secara berurutan kalikan $A$ pertama dan
		ketiga (atau kedua dan keempat) dari
		$A\otimes A\otimes(A\otimes A\otimes M_1\otimes M_2)$ dari kiri ke
		$M_1$ (atau $M_2$);
		\item mula-mula terapkan $\mu:A\otimes A\to A$ pada dua posisi pertama,
		kemudian terapkan $\mu_{M_1\otimes M_2}$. Karena
		% segment-id: o014.li2.chapter7.coalgebra-to-hopf-a.q019
		\[ \Delta\mu = (\mu \otimes \mu) (\identity \otimes c(A, A) \otimes \identity)(\Delta \otimes \Delta), \]
		pemetaan ini sama dengan komposisi
		% segment-id: o014.li2.chapter7.coalgebra-to-hopf-a.q020
		\begin{multline*}
			A \otimes A \otimes M_1 \otimes M_2 \xrightarrow{\Delta \otimes \Delta \otimes \identity_{M_1} \otimes \identity_{M_2}} A \otimes A \otimes A \otimes A \otimes M_1 \otimes M_2 \\
			\xrightarrow{\identity_A \otimes c(A, A) \otimes \identity_A \otimes \identity_{M_1} \otimes \identity_{M_2}} \boxed{A \otimes A \otimes A \otimes A} \otimes M_1 \otimes M_2 \\
			\xrightarrow{\mu \otimes \mu \otimes \identity_{M_1} \otimes \identity_{M_2}} A \otimes A \otimes M_1 \otimes M_2 \\
			\xrightarrow{\identity_A \otimes c(A, M_1) \otimes \identity_{M_2}} A \otimes M_1 \otimes A \otimes M_2 \xrightarrow{\mu_1 \otimes \mu_2} M_1 \otimes M_2.
		\end{multline*}
		Masing-masing perkalian skalar $\mu_1$ dan $\mu_2$ asosiatif dan $c$
		funktorial. Karena itu, komposisi tiga panah terakhir di atas juga dapat
		dipandang sebagai perkalian berurutan $A$ pertama dan kedua (atau
		ketiga dan keempat) dalam bagian berkotak dari kiri ke $M_1$ (atau
		$M_2$), dengan memperhitungkan struktur anyaman.
	\end{enumerate}

	Diagram anyaman membantu memahami hal ini, seperti di bagian akhir
	\cite[\S 7.4]{Li1}. Untuk membuktikan bahwa (a) sama dengan (b), cukup diperiksa
	persamaan berikut\footnote{Ada satu perbedaan kecil yang tidak berbahaya:
	dalam pustaka yang dikutip, morfisme dikomposisikan dari bawah ke atas,
	sedangkan di sini sebaliknya.}
	% segment-id: o014.li2.chapter7.coalgebra-to-hopf-a.g001
	\begin{center}\begin{tikzpicture}
		[baseline=(braid), yscale=0.8]
		\pic[
			braid/number of strands = 6,
			braid/every strand/.style={black},
			name prefix=braid
		] at (0, 0) {braid={s_4 s_2 s_3}};
		\coordinate (braid) at ($(braid-1-s)!0.5!(braid-1-e)$);
		\node[above] at (braid-1-s) {$A$};
		\node[below=0.3em] (P) at (braid-1-e) {$A$};
		\node[above] at (braid-2-s) {$A$};
		\node[below=0.3em] (R) at (braid-2-e) {$A$};
		\node[above] at (braid-3-s) {$A$};
		\node[below=0.3em] at (braid-3-e) {$A$};
		\node[above] at (braid-4-s) {$A$};
		\node[below=0.3em] at (braid-4-e) {$A$};
		\node[above] at (braid-5-s) {$M_1$};
		\node[below=0.3em] (Q) at (braid-5-e) {$M_1$};
		\node[above] at (braid-6-s) {$M_2$};
		\node[below=0.3em] (S) at (braid-6-e) {$M_2$};

		\node (X) [draw, fit=(P) (Q)] {};
		\node[below=0.2em of X] {kalikan pada\; $M_1$};
		\node (Y) [draw, fit=(R) (S)] {};
		\node[below=0.2em of Y] {kalikan pada\; $M_2$};
	\end{tikzpicture} \; = \begin{tikzpicture}
		[baseline=(braid), yscale=0.8]
		\pic[
			braid/number of strands=6,
			braid/every strand/.style = {black},
			name prefix=braid
		] at (0, 0) {braid={s_2 s_4 s_3}};
		\coordinate (braid) at ($(braid-1-s)!0.5!(braid-1-e)$);
		\node[above] at (braid-1-s) {$A$};
		\node[below=0.3em] (P) at (braid-1-e) {$A$};
		\node[above] at (braid-2-s) {$A$};
		\node[below=0.3em] (R) at (braid-2-e) {$A$};
		\node[above] at (braid-3-s) {$A$};
		\node[below=0.3em] at (braid-3-e) {$A$};
		\node[above] at (braid-4-s) {$A$};
		\node[below=0.3em] at (braid-4-e) {$A$};
		\node[above] at (braid-5-s) {$M_1$};
		\node[below=0.3em] (Q) at (braid-5-e) {$M_1$};
		\node[above] at (braid-6-s) {$M_2$};
		\node[below=0.3em] (S) at (braid-6-e) {$M_2$};

		\node (X) [draw, fit=(P) (Q)] {};
		\node[below=0.2em of X] {kalikan pada\; $M_1$};
		\node (Y) [draw, fit=(R) (S)] {};
		\node[below=0.2em of Y] {kalikan pada\; $M_2$};
	\end{tikzpicture}\end{center}
	yang langsung terlihat.

	Untuk menunjukkan bahwa $A\dcate{Mod}$ adalah kategori monoidal, masih
	perlu dibuktikan bahwa kendala asosiativitas dan kendala satuan bagi
	$\otimes$ semuanya merupakan morfisme modul-$A$ kiri. Pernyataan pertama
	tidak sulit. Pernyataan kedua memerlukan sifat koaljabar
	$(\identity\otimes\epsilon)\Delta=\identity_A=
	(\epsilon\otimes\identity)\Delta$, setelah mengidentifikasi
	$A\otimes\munit\simeq A\simeq\munit\otimes A$.
\end{proof}
